% displayname : Diffusion de particules
\documentclass[12pt,fancy]{cours}
\Chapit{2}{4}
\begin{document}
\maketitre{
Dans ce chapitre, nous rencontrons un second exemple de phénomène de transport après la diffusion thermique, celui de la matière dans un milieu immobile: la diffusion de particules. Au niveau microscopique, ce phénomène est causé par un mouvement aléatoire des particules, qui par chocs de proche en proche finissent par homogénéiser le milieu. La diffusion se caractérise par une seule grandeur scalaire, le coefficient de diffusion.
}
{
  \item Cours Thermodynamique 1 - 0
  \item TD Thermodynamique 1
}


%%%%%%%%%%%%%%%%%%%%%%%%%%%%%%%%%%%%%%%%%%%%%%%%%%%%%%%%%%%%%%%%%%%%%%%%%%
\section{Loi fondamentale du transport de particules}

\subsection{Introduction au transport de particules}

\parag{Mécanismes de transport de particules}

Dans le chapitre précédent, nous avons étudié le phénomène de conduction thermique. La conduction correspond au transfert d’énergie sans déplacement macroscopique de matière dans le milieu, et s’effectue par diffusion de l’énergie de proche en proche. L’énergie peut aussi être transportée avec déplacement de matière dans le milieu : il s’agit de convection. Le transport de particules (atomes, molécules, grains, poussière, etc.) dans un milieu peut également se faire selon ces deux modes.
\begin{liste}
\item Transport par \imp{diffusion} : transport de particules de proche en proche dans un milieu immobile.
\item Transport par \imp{convection} : transport de particules avec déplacement macroscopique de la matière constituant le milieu (qui est alors fluide).
\end{liste}

\parag{Mouvement brownien}
À l’échelle microscopique, le transfert d’énergie se fait au gré des chocs entre particules. Au cours de ces chocs, les particules rebondissent dans des directions aléatoires, et peuvent ainsi se déplacer, lentement, de proche en proche, sur des distances mésoscopiques (comme un parfum qui diffuse dans une pièce).
Ce mouvement aléatoire est appelé mouvement brownien.
%La conduction thermique, qui correspond à un transport d’énergie interne, nécessite donc également un transport des particules.
\begin{definition}
On appelle \imp{mouvement brownien}
 le mouvement désordonné d’une particule perturbé par les chocs incessants des particules voisines.
\end{definition}
Ces chocs sont d’autant plus violents que l’\imp{agitation thermique} est élevée. Ce mouvement a été mis en évidence par le botaniste Robert Brown sur des grains de diamètre d’environ $\SI{1}{\mu m}$, agités par les chocs avec les molécules d’eau ; le mouvement brownien est donc même visible à l’échelle mésoscopique. Ces expériences ont été reproduites au début du \textsc{xx}$^{\rm e}$ siècle par Jean Perrin (\Fig{perrin}).

\begin{figure}[h!]
\centering
\includegraphics[scale=0.4]{perrin.png}
\caption{Reproduction d'un dessin de Jean Perrin dans \emph{Mouvement brownien et réalité moléculaire}. Sont représentées ici trois trajectoires de particules d'environ 1 µm de diamètre, dont la position est repérée toutes les 30s. Ces trajectoires montrent l'existence de chocs à l'échelle microscopique avec les molécules du fluide environnant. \Wiki}
\label{fig:perrin}
\end{figure}
%%%%%%%%%%%%%

Pour décrire le transport de particules, nous procéderons comme pour la conduction thermique.
\begin{liste}
\item Nous allons définir le vecteur associé au courant de la grandeur transportée (ici le nombre de particules $N$, voir \Sec{def}).
\item Nous allons établir une équation de conservation (ici la conservation de la matière, voir \Sec{conserv}).
\item Nous allons proposer une loi phénoménologique pour relier l’inhomogénéité (ici de la concentration en particules) au courant de matière, voir \Sec{fick}. En combinant avec l’équation de conservation, nous aboutirons alors à une équation de diffusion, voir \Sec{equadiff}.
\end{liste}
%. N, comme pour la conduction thermique. Nous discuterons plus en détails des propriétés de cette équation, et nous proposerons plusieurs modèles microscopiques.

\subsection{Flux et densité de courant particulaires}
\label{sec:def}
\parag{Densité particulaire} Nous faisons l’hypothèse de l’équilibre thermodynamique local (ETL) : nous supposons que tout volume élémentaire $\d V$ du milieu centré en un point $M$, de taille mésoscopique $\delta \ll L$ petite devant la taille $L$ du milieu, est à l’équilibre thermodynamique. Les variables thermodynamiques de ce volume, température $T(M,t)$, pression $p(M,t)$, etc. sont donc bien définies.
Soit $\d N(M,t)$ le nombre de particules dans ce volume élémentaire $\d V$.
\begin{definition}
On appelle \imp{densité particulaire} $n(M,t)$ le nombre de particules par unité de volume, soit
$$
\d N(M,t) = n(M,t) \d V.
$$
\end{definition}

%Nous nous plaçons dans l’hypothèse de l’équilibre thermodynamique local.
\noindent
La densité particulaire $n(M,t)$ se mesure $\SI{}{m^{-3}}$.

%%%%%%%%%%%%%%%%%%%%%%%%
\begin{figure}[b!]
\centering
\subfigure[]{\label{fig:jN}
\includegraphics[scale=1.1]{def_flux_N.pdf}
}\hspace{1cm}
\subfigure[]{\label{fig:bilan_N}
\includegraphics[width=0.45\textwidth]{bilan_jN.pdf}}
\caption{(a) Vecteur densité de courant de particules $\Vec{\jN}$. (b) Bilan de particules à une dimension.}
\end{figure}
%%%%%%%%%%%%%%%%%%%%%%%%

\parag{Densité de courant et flux de particules}
Soit une section $\Surf$ séparant deux milieux $1$ et $2$ (ou deux portions d’un même milieu), de normale locale $\Vec{n}_{1\rightarrow2}$ dirigée de $1$ vers $2$. Nous définissons le flux de particules et la densité de courant de particules de la même manière que le flux et la densité de courant thermique.
On observe que le nombre de particules élémentaire traversant cette section du milieu $1$ vers le milieu $2$, noté $\delta N_{1\rightarrow 2} $, est proportionnel à l’aire $\d S$ de la section et à la durée $\d t$ des échanges (\Fig{jN}). %Nous pouvons donc mettre $\delta Q$ sous la forme d’un flux élémentaire.
\begin{definition}
On appelle \imp{flux de particules} $\Phi_N$ (en \SI{}{s^{-1}}) le nombre de particules traversant une section donnée par unité de temps, soit
\begin{equation*}\label{eq:jth}
\Phi_N = \frac{\delta N_{1\rightarrow 2}}{\d t}.
\end{equation*}
% \begin{equation*}\label{eq:jth}
% \Phi_N = \dd{ N_{1\rightarrow 2}}{t} = \iint_{\mc{S}}\Vec{\jN}\cdot\vv{\d S}.
% \end{equation*}
\end{definition}
\begin{remarque}
Attention à ne pas confondre le flux de particules avec un flux thermique ! Par simplicité, nous notons souvent les deux quantités $\Phi$, mais elles n’ont pas la même dimension.
\end{remarque}
\begin{definition}\label{def:jth}
On définit le vecteur \imp{densité de courant de particules} $\Vec{\jN}$ (en \SI{}{s^{-1}.m^{-2}}) tel que
\begin{equation*}
\Phi_N = \iint_{\Surf}\Vec{\jN}\cdot\vv{\d S}.
\end{equation*}
\end{definition}
\noindent
%Rappelons que le vecteur $\Vec{\dS}$ est de norme la surface élémentaire $\d S$ et est porté par le vecteur unitaire $\Vec{n}_{1\rightarrow2}$.
% de $1$ vers $2$.
 % La densité de courant thermique $\jN$ est homogène à l’inverse d’un temps et d’une surface (en $\SI{}{s^{-1}.m^{-2}}$).
La direction de $\Vec{\jN}$ indique la direction moyenne de déplacement des particules.

\subsection{Bilan de matière}
\label{sec:conserv}

\begin{methode}[Établir l'équation de conservation de la matière]
  \item Définir le système comme une portion élémentaire du milieu :
  \begin{liste}
    \item tronçon d'épaisseur $\d x$ et de section $S$ en cartésienne ;
    \item coquille cylindrique de rayon $r$, d'épaisseur $\d r$ et de hauteur $H$ en cylindrique ;
    \item coquille sphérique de rayon $r$, d'épaisseur $\d r$ en sphérique.
  \end{liste}
  \item Exprimer la variation d'énergie interne $\d N$ pendant $\d t$ en fonction de $\dr{n}{t}$.
  \item Exprimer cette variation $\d N$ en fonction des dérivées spatiales de $\Vec{\jth}$.
  \item Simplifier et en déduire l'équation de conservation.
\end{methode}


\parag{Bilan à une dimension} Réalisons un bilan de matière sur une portion de conduite homogène de section $S$ et de longueur $\d x$, donc de volume $\d V = S\d x$. La densité particulaire  $n(x,t)$ est supposée uniforme sur la section, et donc ne dépend que de $x$. Le mouvement des particules est limité à la seule direction $(Ox)$, et donc $\Vec{\jN} = \jN(x) \ex$ est porté par $\ex$, comme illustré sur la \Fig{bilan_N}.

%\parag{Expression du nombre de particules}
La variation $\d N$ du nombre de particules dans la conduite entre $t$ et $t+\d t$ correspond par définition au nombre de particules présentes à l'instant $t+\d t$, moins le nombre qui était présent à l'instant $t$, autrement dit
\begin{equation}
\label{eq:etablissement_eq_diffusion_bilan1}
\d N = \underbrace{n(x,t+\d t)\d V}_{\text{nb particules à } t + \d t} -\underbrace{n(x,t)\d V}_{\text{nb particules à } t} =n(x,t+\d t)S\d x-n(x,t)S\d x=\frac{\partial n}{\partial t}(x,t)\,\d t\,S\d x,
\end{equation}
par développement de Taylor à l’ordre $1$.
Puisque $n$ représente une densité volumique de particules, on multiplie bien par le volume élémentaire $\d V$ pour obtenir un nombre de particules.

Or, la variation $\d N$ correspond au nombre de particules $\delta N_e$ qui sont rentrés pendant $\d t$ (par la section d'abscisse $x$), moins le nombre $\delta N_s$ de celle qui sont sorties pendant $\d t$ (par la surface d'abscisse $x+\d x$):
\begin{align}
\label{eq:etablissement_eq_diffusion_bilan2}
\nonumber
\d N = \delta N_e - \delta N_s = \Phi(x,t)d t-\Phi(x+\d x,t)  \d t
&=\jN(x,t) S \d t- \jN(x+\d x,t) S\d t\\
& =-\frac{\partial \jN}{\partial x}(x,t) \d x S\d t,
\end{align}
toujours d’après un développement de Taylor à l’ordre $1$.
%où nous avons noté $j_x$ la composante selon $\vv{e_x}$ du courant $\vv{j}$.

\parag{Équation de conservation de la matière}
En égalisant les Éq.~\eqref{eq:etablissement_eq_diffusion_bilan1} et~\eqref{eq:etablissement_eq_diffusion_bilan2} on obtient
\begin{equation}
\label{eq:equation_continuite_1D}
\dfrac{\partial n}{\partial t}+\dfrac{\partial \jN}{\partial x}=0.
\end{equation}
On peut réaliser le même type de bilan à trois dimensions, ce qui fait apparaître l’opérateur divergence.

\begin{theorem}[Équation de conservation de la matière]\label{form:equation_continuite_3D}
\begin{equation*}
\dfrac{\partial n}{\partial t}+\dive{\Vec{\jN}}=0.
\end{equation*}
\end{theorem}
\noindent On appelle également cette équation, \imp{équation de continuité}.

%On peut réécrire de façon plus compacte cette équation (et la rendre également indépendante du système de coordonnées) si on reconnaît l'expression de l'opérateur vectoriel divergence

\begin{remarque}
Dans la démonstration précédente, nous avons considéré qu'il ne pouvait pas y avoir de création spontanée de particules au sein du milieu. Il existe des cas où cette hypothèse est fausse: c'est par exemple le cas d’un bilan du nombre de neutrons dans un réacteur nucléaire, car la réaction de fission produit des neutrons. Pour rendre compte de ces phénomènes, il suffit d’ajouter un terme source $p_{\rm vol}$ (en $\SI{}{s^{-1}.m^{-3}}$) dans le second membre de la \Formu{equation_continuite_3D}, ce qui donne
\begin{equation}
\dfrac{\partial n}{\partial t}+\dive{\vv{j}}= p_{\rm vol}.
\end{equation}
 Ce terme est positif si les particules sont créées, négatif si elles sont détruites.
 \end{remarque}

\section{Équation de diffusion}

\subsection{Loi de Fick}
\label{sec:fick}

%\parag{Loi de Fick}
La \Formu{equation_continuite_3D} est une loi fondamentale : elle traduit la conservation de la matière. Malheureusement, elle fournit une seule équation pour deux inconnues (la densité $n$ et le vecteur densité de courant thermique $\Vec{\jN}$). Pour résoudre cette équation et déterminer la densité, il faut relier $\Vec{\jN}$ et $n$. C'est le rôle de la loi phénoménologique de \imp{Fick}.
%, qui fournit une expression de $\Vec{\jN}$ valable pour les phénomènes de conduction.

\begin{theorem}[Loi de Fick]\label{form:fick}
\begin{equation*}
\Vec{\jN}=-D\,\Grad{n}.
\end{equation*}
La grandeur $D$, appelée \imp{coefficient de diffusion}, se mesure en $\SI{}{\meter^2.s^{-1}}$.
\end{theorem}
\noindent  On distingue trois types de diffusion selon l’état physique du milieu diffusant.
%%%%%%%%%%%%%%%%%%%%%%
\begin{figure}[H]
\centering
% \includegraphics[width=0.8\textwidth]{D_ODG.png}

\begin{tabular}{|c|c|c|c|c|}
\hline Phase & Milieu de diffusion & Particules diffusantes & Température $\left({ }^{\circ} \mathrm{C}\right)$ & $D\left(\mathrm{m}^2 \cdot \mathrm{s}^{-1}\right)$ \\
\hline \multirow{3}{*}{$\begin{array}{c}\text { Gaz (pression } \\
\text { atmosphérique) }\end{array}$} & \multirow{3}{*}{ Air } & \multirow{2}{*}{ Dihydrogène } & 25 & $6,3.10^{-5}$ \\
\cline { 4 - 5 } & & & 400 & $3,24.10^{-4}$ \\
\cline { 3 - 5 } & & \multirow{2}{*}{ Vapeur d'eau } & 25 & $2,4.10^{-5}$ \\
\cline { 4 - 5 } & & & 400 & $1,13.10^{-4}$ \\
\hline \multirow{3}{*}{ Liquide } & \multirow{2}{*}{ Eau } & Sucre & 25 & $5,2.10^{-10}$ \\
\cline { 3 - 5 } & & Éthanol & 25 & $1,2.10^{-9}$ \\
\cline { 2 - 5 } & Benzène & Éthanol & 25 & $3,0.10^{-9}$ \\
\cline { 2 - 5 } & Argent fondu & Or & 980 & $2,5.10^{-9}$ \\
\hline Solide & Cuivre solide & Aluminium solide & 20 & $1,3.10^{-30}$ \\
\hline
\end{tabular}
\caption{Quelques valeurs typiques de coefficients de diffusion.
% à température ambiante, pour un gaz (l’air), un liquide (l’eau) et un solide (le cuivre).
}
% \label{fig:}
\end{figure}
%%%%%%%%%%%%%
\begin{liste}
\item La diffusion de particules dans un gaz. Les particules peuvent être des molécules, mais aussi des particules plus grosses (poussières, etc.). Le coefficient de diffusion vaut typiquement \eqbox{D_{\rm gaz} \sim \SI{e-4}{}.}
% pour de la diffusion de molécules dans l’air.
\item La diffusion de particules dans une solution liquide. Les particules peuvent être des molécules du solvant, des ions, ou des solutés. Le coefficient de diffusion vaut typiquement \eqbox{D_{\rm liq} \sim \SI{e-9}{}.}
 % pour de la diffusion en solution.
\item La diffusion de particules dans un solide cristallin (d’impuretés dans les semi-conducteurs par exemple). Le coefficient de diffusion dans les solides peut être aussi faible que \eqbox{D_{\rm sol} \sim \SI{e-30}{}.}
\end{liste}

  La diffusion est la plus rapide dans les gaz, elle est plus lente dans les liquides, et extrêmement lente dans les solides.

La loi de Fick, comme la loi de Fourier, est une loi phénoménologique : elle permet d’expliquer les observations sans proposer de mécanisme sous-jacent. On peut la justifier par les mêmes arguments que pour la loi de Fourier (faibles variations de densité).


\subsection{Équation de diffusion particulaire}
\label{sec:equadiff}
Nous souhaitons maintenant obtenir une équation générale portant sur la densité particulaire $n(M,t)$. Il suffit de
combiner la loi de Fick \Formu{fick} avec la loi de conservation de la matière en l’absence de terme
source \Formu{equation_continuite_3D}. On obtient alors l’équation de la diffusion,
\begin{theorem}[Équation de la diffusion particulaire]\label{form:eqdiff}
\begin{equation*}
\dfrac{\partial n}{\partial t}=D\,\Delta n.
\end{equation*}
\end{theorem}

%Cette équation généralise à 3 dimensions en\footnote{La divergence d'un gradient s'identifie au Laplacien scalaire, voir le chapitre de rappels mathématiques (analyse vectorielle).}
%\begin{equation}
%\begin{split}
%\dfrac{\partial n(M,t)}{\partial t}-D\,\Delta n(M,t)=0.
%\end{split}
%\end{equation}
Il s’agit de la même équation que pour la conduction thermique (seule la constante de diffusion change puisque $D$ remplace ici la diffusivité thermique). L’équation de la \Formu{eqdiff} est une \imp{équation différentielle aux dérivées partielles, linéaire, et du second ordre} pour la densité de particules $n$. Pour résoudre cette équation, il est possible d'utiliser le formalisme complexe, comme pour la diffusion thermique et l’effet de cave
% ou de façon équivalente l'analyse de Fourier (décomposition d'un signal quelconque en harmoniques, étude de la réponse à une harmonique quelconque puis somme de toutes les réponses individuelles pour reconstituer la réponse globale)
,
ou encore de chercher une solution auto-similaire (\Sec{regime_autosimilaire}). Il faut adjoindre à cette équation deux conditions aux limites ($\Delta = 2^{\rm nd}$ ordre en l’espace)  et une condition initiale ($\dr{}{t} = 1^{\rm er}$ ordre en temps). On distingue deux types de conditions aux limites.
\begin{liste}
\item Soit la densité de particules $n$ est imposée.
\item Soit le flux de particules $\Phi_N$ est imposé.
\end{liste}
%Il peut aussi arriver qu’une condition soit mixte.

Aucune connaissance sur la solution générale d'une équation de diffusion n'est au programme, et vous serez donc (normalement) guidés dans les problèmes (voir TD).
%Néanmoins, si vous souhaitez savoir d'où vient la forme de la solution générale, vous pouvez vous reporter au Complément~\ref{comp:solution_générale_eq_diff}.
En revanche, les cas très simples (régime stationnaire en particulier) sont fréquents dans les problèmes et il faut maîtriser l'obtention de la solution.


%\begin{remarque}
%S'il y a un terme source volumique $p(M,t)$, il faut toujours en tenir compte dans le membre de droite, car celui-ci intervient dans le flux de particules.
%\end{remarque}

\subsection{Propriétés de l'équation de diffusion}

\parag{Loi d'échelle}
Étudions un processus diffusif de coefficient de diffusion $D$. Par exemple, un traceur coloré est injecté en un point  d’un liquide. Soit $t$ le temps écoulé depuis l’injection du traceur, et $\delta$ la distance caractéristique sur laquelle le traceur s’est propagé. En ordre de grandeur, l'équation de diffusion se réécrit
\begin{equation}
\begin{split}
\dfrac{\partial n}{\partial t}=D\Delta n \quad\Leftrightarrow\quad\dfrac{n}{t}\sim D\dfrac{n}{\delta^{2}} \quad\Leftrightarrow\quad \delta \sim \sqrt{Dt}.
\end{split}
\end{equation}

\begin{prop}\label{prop:delta}
Les particules diffusées se propagent sur une distance $\delta$ proportionnelle à la racine du temps $t$ suivant
$$\delta \sim \sqrt{Dt}.$$
\end{prop}
\noindent La diffusion est un phénomène lent, contrairement à un processus balistique (où un objet se propage à vitesse constante), pour lequel la distance est proportionnelle au temps.
% pour le remarquer, définissons les longueurs et temps caractéristiques du phénomène
%\begin{equation}
%\begin{split}
%\delta=\sqrt{D\tau}\qquad\text{et}\qquad \tau=\frac{\delta^2}{D}.
%\end{split}
%\end{equation}
Pour donner quelques ordres de grandeur, la diffusion des molécules dans un gaz, de coefficient de diffusion $D\sim\SI{e-5}{m^2.s^{-1}}$, sur une distance de l'ordre de $\SI{1}{m}$, prend de l'ordre de $\SI{e5}{s}$, soit environ un jour ! Comme la diffusion est un phénomène lent, elle est généralement négligeable dans les fluides devant la convection.
%\begin{remarque} Le temps et la diffusion jouent exactement le même rôle ! Application : vieillissement artificiel d'un objet. \end{remarque}

\medskip


\parag{Irréversibilité}
L'équation de diffusion fait intervenir une \imp{dérivée première en temps}. Par conséquent, elle n'est pas invariante par renversement du temps $t\rightarrow -t$ (nous verrons que l’équation de propagation des ondes, qui fait intervenir des dérivées secondes par rapport au temps, est invariante par renversement du temps).
\begin{prop}\label{prop:}
  La diffusion est un phénomène \imp{irréversible}.
\end{prop}
 Cela est conforme à l'intuition: si on lâche une goutte d'encre dans un verre d'eau, l'encre s'étale dans tout le verre jusqu'à homogénéisation complète. Mais le processus inverse, où en partant d'encre diluée dans l'eau on rejoindrait un état final où la goutte d'encre se reforme au centre, ne s'observe pas.
%Nous retrouverons ce phénomène irréversible dans le chapitre suivant sur la diffusion thermique, qui se traduit par la création d'entropie.
Notons enfin que l'irréversibilité se voit aussi directement dans la loi de Fick, car si on renverse la flèche du temps, on inverse le vecteur courant $\vv{j}$, mais ni le gradient de densité ni le coefficient de diffusion ne changent de signe.
%" La loi de Fick n'est pas non plus invariante par renversement du temps et traduit un phénomène irréversible.

\subsection{Régime stationnaire}

Par analogie avec la diffusion thermique, le flux de particules est conservé en régime stationnaire. Cela se traduit par les lois suivantes (l’une est locale, l’autre est globale ou intégrale).


\begin{prop}[Conservation du flux de particules -- forme locale]\label{prop:conservj}
En régime stationnaire,
\begin{equation*}
\Dive{\Vec{{\jN}}}=0.
\end{equation*}
\end{prop}
%Dans notre exemple à une dimension, $\Vec{\jth} = \jth(x) \ex$ et donc $\d \jth/\d x = 0$ : $\jth(x)$ est constant. Comme $\jth$ est aussi uniforme sur une section, on en déduit que $\Phi(x) = \jth S$ est constant également : le flux thermique $\Phi$ ne dépend pas de $x$, et donc de la section $\mc{S}$ sur laquelle on le calcule. Généralisons ce résultat.

 On appelle \imp{tube de champ} du vecteur densité de courant de particules $\Vec{\jN}$ un ensemble de lignes tangentes à $\Vec{\jN}$ qui s'appuient sur deux sections $\S_1$ (en entrée) et $\S_2$ (en sortie).

\begin{figure}[h!]
\centering
\includegraphics[width=0.5\textwidth]{bilan_jN_tube.pdf}
\caption{Tube de champ du vecteur densité de courant de particules $\Vec{\jN}$.}
\label{fig:tube}
\end{figure}
%%%%%%%%%%%%%

\begin{prop}[Conservation du flux particules -- forme globale]
En régime stationnaire,
$\Phi_N = \cte$
 le long d'un tube de champ ($\Phi_N$ est indépendant de la section $\Surf$).
\end{prop}
\noindent On dit que $\Vec{\jN}$ est à \imp{flux conservatif}.

\begin{exemple}{Diffusion de particules dans une conduite}
Une conduite de section circulaire (surface $S$) et d’axe $(Ox)$ est remplie d’eau. En entrée $x=0$, de l’éthanol est injecté en continue de façon à maintenir une densité de particules constante $n_1$. La densité en éthanol vaut $n_2$ en $x=L$.
On suppose le régime permanent atteint. Déterminer le profil de densité  en éthanol dans l’eau.
\begin{quest}
  \item De quelles variables dépend $n$ ? Comment est orientée le vecteur $\Vec{\jN}$ ?
  \item Écrire et résoudre l'équation de diffusion en régime permanent.
  \item Utiliser les conditions aux limites pour déterminer complètement $n$.
  \item Calculer le flux de particules $\Phi_N$ et commenter.
\end{quest}
\end{exemple}
\examplesol{
\begin{quest}
\item Géométrie : la densité particulaire est uniforme sur une section, et indépendante du temps : $n(x)$. On en déduit par la loi de Fick que la densité de courant de particules vaut $\Vec{\jN} = - D\, \grad n = - D  \dd{n}{x}\ex$.
\item Équation de diffusion : en régime stationnaire, on a $\dive{\Vec{\jN}} = 0$ et donc
\eqbox{ \Delta n = \ddm{n}{x}{2} = 0.
}
La solution est de la forme $n(x) = A+ Bx$.
\item Conditions aux limites : $n(x=0) = A = n_1$ puis $n(x=L) = A+ BL = n_2$ donc $B = (n_2-n_1)/L$. On en déduit le profil de densité particulaire
\eqbox{
n(x) = (n_2-n_1)\frac{x}{L} + n_1.
}
\item Remarque : le flux de particule vaut $\Phi_N = - D \dd{n}{x} S = \frac{DS}{L}(n_1-n_2)$, soit
\eqbox{\Phi_N = \frac{DS}{L}(n_1-n_2).}
 Il est bien indépendant de $x$ et donc conservé le long de la conduite, comme attendu en régime stationnaire.
\end{quest}
}

\section{Modèles microscopiques}

\begin{definition}[Marche aléatoire]
Une \imp{marche aléatoire} est un modèle mathématique d'un système dont la dynamique est décrite par une succession de pas \imp{aléatoires indépendants}.
\end{definition}

La marche aléatoire est un modèle simple qui permet d'interpréter physiquement le phénomène de diffusion de particules à l'échelle microscopique : le mouvement brownien. 

\subsection{Marche aléatoire : modèle probabiliste}
\label{sec:marche}

\parag{Modèle discret unidimensionnel isotrope}
Pour interpréter physiquement le phénomène de diffusion de particules à l'échelle microscopique, considérons le modèle de marche aléatoire suivant :
\begin{enumerate}
  \item \udl{unidimensionnel} : la particule se déplace le long d'un axe $(Ox)$. Elle part de la position initiale $x=0$.
  \item \udl{discret} : la particule effectue des sauts de longueur constante $\ell$ et de durée constante $\tau$. Elle ne peut donc occuper que les sites d'un réseau de positions $x_n = n \ell$ aux instants $t_m = m \tau$, avec $n,m \in \mathbb{N}$.
  \item \udl{isotrope} : à chaque saut, la particule a la même probabilité de se déplacer vers la droite ou vers la gauche : ces probabilités valent donc $1/2$.
\end{enumerate}

\begin{figure}[H]
  \centering
  \includegraphics[scale=1.2]{random_walk_1D.pdf}
  \caption{(a) Modèle de marche aléatoire unidimensionnelle. (b) Relation de récurrence.}
  \label{fig:}
\end{figure}

\parag{Propriétés statistiques}
On décrit les propriétés statistiques de la marche aléatoire par la fonction de probabilité $P(x,t)$, définie comme la probabilité que la particule soit en $x$ à l'instant $t$. Cette fonction est normalisée, i.e.  la particule est nécessairement quelque part,
\begin{equation}
\sum_{n \in \Z} P(x_n,t) = 1, \quad \forall t.
\end{equation}
Si la particule est à l'instant $t_{m+1}$ en $x_{n}$, c'est qu'à l'instant  $t_{m}$ précédent, elle était ou bien en $x_{n-1}$ ou bien en $x_{n+1}$, avec la même probabilité. La fonction $P$ vérifie donc la relation de récurrence
\eqbox[p(xn,t)]{
  P(x_n,t_{m+1})=\frac{1}{2}\,P(x_{n-1},t_m)+\frac{1}{2}\,P(x_{n+1},t_m).
  }
Les propriétés statistiques de la marche aléatoire sont donc entièrement déterminées par la relation de récurrence ci-dessus, et par la condition initiale $P(x,0) = 0$  si $x \neq 0$ et $P(0,0) = 1$. On peut déterminer :
\begin{liste}
  \item La \udl{position moyenne} :
  \begin{equation}
    \moy{x}\!(t) = \sum_{n \in \Z} x_n \, P(x_n,t).
  \end{equation}
  Comme la particule a la même probabilité de sauter à droite ou à gauche, la fonction $P(x,t)$ est paire en $x$ pour tout $t$. Par conséquent, la position moyenne vaut 
  \eqbox{
    \moy{x}\!(t) = 0.
    }
    \item La \udl{distance quadratique moyenne} :
    \begin{equation}
      \moy{x^2}\!(t) = \sum_{n \in \Z} x_n^2 \, P(x_n,t).
    \end{equation}
    On peut montrer que la distance quadratique moyenne vérifie  
    \eqbox{
      \moy{x^2}\!(t) = \frac{t\ell^2}{\tau}.
    }
    \begin{demo}
    Nous avons la relation de récurrence
    \begin{align}
      \moy{x^2}\!(t_{m + 1}) = \sum_{n \in \Z} x_n^2 \, P(x_n,t_{m + 1}) & = \sum_{n \in \Z} x_n^2 \left[ \frac{1}{2} P(x_{n-1},t_m) + \frac{1}{2} P(x_{n+1},t_m) \right]. \\ 
      % \intertext{En changeant l'indice de sommation $n$ en $n+1$ dans le premier terme, et en $n-1$ dans le second, on trouve}
      &  = \sum
      _{n \in \Z} \frac{1}{2}(x_{n+1}^2 + x_{n-1}^2)  P(x_n,t_m) \\ 
      & = \sum_{n \in \Z} (x_n^2 + \ell^2) P(x_n,t_m). \\
      % \intertext{On en déduit la relation de récurrence}
      & = \moy{x^2}\!(t_m) + \ell^2.
    \end{align} 
    On en déduit la relation annoncée
    \begin{equation}
      \moy{x^2}\!(t_m) = m \ell^2 = \frac{t_m \ell^2}{\tau}.
    \end{equation}
  \end{demo}
    % \item La \udl{distance typique} :
    % \begin{equation}
    %   \delta(t) = \sqrt{\moy{x^2}\!(t)}.
    % \end{equation}
    % Elle vérifie  la loi caractéristique de la diffusion,
    % \eqbox{
    % \delta(t) \propto \sqrt{t}.
    % }
\end{liste}
\parag{Lien avec la diffusion} La marche aléatoire permet de modéliser le phénomène de diffusion à l'échelle microscopique. Définissons la distance typique parcourue par la particule après un temps $t$,
\begin{equation}
  \delta(t) = \sqrt{\moy{x^2}\!(t)}.
\end{equation}
D'après les propriétés précédentes, elle vérifie la loi caractéristique de la diffusion,
\eqbox{
\delta(t) \propto \sqrt{t}.
}
De plus, déterminons l'équation différentielle satisfaite par la fonction de probabilité $P(x,t)$. On se place dans l'approximation du continu : on considère que $P$ est une fonction continue de $x$ et $t$, et que $P(x,t)$ varie très lentement ser les échelles spatiale $\ell$ et temporelle $\tau$.  La relation de récurrence s'écrit alors 
\begin{equation}
\label{eq:p(xn,t)}
P(x_n,t+\tau)=\frac{1}{2}P(x -\ell, t)+\frac{1}{2}P(x + \ell,t).
\end{equation}
Simplifions par des développements limités.
\begin{liste}
  \item \udl{DL$_2(x) :$} 
  \begin{equation}
  \begin{split}
  P(x \pm \ell,t)&= P(x+\ell,t)\simeq P(x,t)\pm \ell\frac{\partial P}{\partial x}(x,t)+\frac{\ell^2}{2}\frac{\partial^2 P}{\partial x^2}(x,t).\\
  \end{split}
  \end{equation}
  \item \udl{DL$_1(t) :$}
  \begin{equation}
  P(x,t+\tau)\simeq P(x,t)+\tau\frac{\partial P}{\partial t}(x,t).
  \end{equation}
\end{liste}
On obtient alors 
\begin{equation}
  P(x,t)+\tau\frac{\partial P}{\partial t}(x,t) = P(x,t) + \frac{\ell^2}{2}\frac{\partial^2 P}{\partial x^2}(x,t),
\end{equation}
soit l'équation de diffusion 
% \begin{equation}
\eqbox{
  \frac{\partial P}{\partial t}= \frac{\ell^2}{2\tau} \frac{\partial^2 P}{\partial x^2},
  }
% \end{equation}
avec pour coefficient de diffusion $D=\frac{\ell^2}{2\tau}$.  Pour une marche aléatoire en $1$D, on a donc $\delta = \sqrt{2Dt}$.


% Le facteur $1/2$ est spécifique à la dimension du problème ($1D$), mais en ordre de grandeur on aura toujours
% \begin{equation}
%   D \simeq \frac{\ell^2}{\tau}.
% \end{equation}
% Ce modèle met en évidence que la diffusion des particules peut s'interpréter par des chocs ou sauts aléatoires. Grâce à ce modèle, nous pouvons définir précisément la distance caractéristique de diffusion $\delta$ introduite dans la \Prop{delta}. On pourrait choisir naïvement la position moyenne $\moy{x} = \sum_x x \, p(x,t)$, mais comme la particule a la même chance de sauter à droite ou à gauche, la position moyenne vaut $\moy{x} = 0$ ! Par contre, on peut choisir la distance quadratique moyenne
% $\delta = \sqrt{\moy{x^2}}$ avec $\moy{x^2} = \sum_x x^2 \, p(x,t)$. D’après la \Prop{delta}, on s’attend donc à ce que $\sqrt{\moy{x^2}} \sim \sqrt{D t}$ dans un processus diffusif.

% En supposant que $p$ varie lentement sur une distance de l’ordre de $a$, on peut effectuer un développement limité à l’ordre $2$,

% De même, en supposant que $p$ varie lentement sur un temps de l’ordre de $\tau$, on a
% \begin{equation}
% p(x_n,t+\tau)\simeq p(x_n,t)+\tau\,\frac{\partial p}{\partial t}(x_n,t).
% \end{equation}
% En utilisant ces développements limités dans l'Éq.~\eqref{eq:p(xn,t)} on remarque que $p(x_n,t)$ obéit à l’équation
% \begin{equation}
% \frac{\partial p}{\partial t}=\frac{\ell^2}{2\tau}\frac{\partial^2 p}{\partial x^2}.
% \end{equation}


\parag{Simulation numérique}

On se reportera à la \refCN{3}, pour l'étude de la diffusion d'un grand nombre de particules. Voici un modèle de marche aléatoire $1$D en \texttt{Python} pour une seule particule. Nous commençons par définir les grandeurs $\ell$, $\tau$ et $D$, puis nous simulons \texttt{Nstep} sauts de la particule. La fonction \texttt{marche} renvoie la distance au carré parcourue après \texttt{Nstep} sauts.
%Le code ci-dessous représente une marche aléatoire à deux dimensions.

\inputminted{python}{Code/marche1.py}
% \python{marche1}

On ne connaît pas la fonction $p(x,t)$, mais on peut estimer la distance quadratique moyenne
$\delta = \sqrt{\moy{x^2}}$ en moyennant sur un grand nombre de simulations. On peut ensuite calculer la distance quadratique moyenne.
\inputminted{python}{Code/marche2.py}
% \python{marche2}

On peut ensuite calculer $\delta(t)$ pour différents temps de simulations $t$. On remarque graphiquement que la loi $\delta(t) \sim \sqrt{Dt}$ est bien vérifiée (\Fig{marche_1D}). On peut d’ailleurs montrer mathématiquement la relation exacte
\begin{equation}
\sqrt{\moy{x^2}} = \sqrt{2Dt}.
\end{equation}
Le facteur $2$ est valable à une dimension.
\inputminted{python}{Code/marche3.py}
% \python{marche3}



\begin{figure}[h!]
\centering
\subfigure[]{\label{fig:marche_1D_simulation}
\includegraphics[scale=1]{marche1D_traj.pdf}
}
\subfigure[]{\label{fig:marche_1D_loi_diffusion}
\includegraphics[scale=1]{marche1D_loi_diffusion.pdf}
}
% \subfigure[]{\label{fig:marche_2D}
% \includegraphics[scale=0.5]{rand_walk2000.png}}
\caption{(a) Trois simulations 1D d'une marche aléatoire avec $N=1000$ sauts. On représente la position $x$ (en unité de $\ell$) en fonction du temps $t$ (en unité de $\tau$). (b) Distance typique $\delta = \sqrt{\moy{x^2}}$ en fonction du temps~$t$. La loi caractéristique de la diffusion 1D, $\delta = \sqrt{2Dt}$ (traits rouges) est bien vérifiée.}
% \label{fig:}
\end{figure}
%%%%%%%%%%%%%

\begin{figure}[H]
  \centering
  \subfigure[]{\label{fig:marche_2D}
  \includegraphics[height=6.5cm]{rand_walk5000.pdf}
  }
  \subfigure[]{\label{fig:perrin}
  \includegraphics[height=6.5cm]{perrin.png} 
  }
  \caption{(a) Simulation 2D d'une marche aléatoire avec $N=5000$ sauts. On représente la trajectoire dans le plan $(x,y)$ de la particule. (b) Trajectoire expérimentale d'une particule de pollen dans l'eau, observée par Perrin en 1908, caractéristique du mouvement brownien.}
  \label{fig:marche_2D}
\end{figure}

\subsection{Marche aléatoire : modèle du gaz parfait}

Plaçons nous dans le cas d'un gaz, supposé parfait, de masse molaire $M$. Nous voulons estimer le coefficient de diffusion d’une molécule de ce gaz dans le gaz lui-même. Comme le gaz est parfait, les chocs entre molécules sont instantanés et les molécules ne sont soumises à aucune force entre deux chocs : elles se propagent alors en ligne droite à vitesse constante $v$. Nous pouvons donc décrire un gaz parfait avec un modèle de marche aléatoire à trois dimensions pour un gaz parfait.

La distance $\ell$ introduite dans le modèle de marche aléatoire \Sec{marche} correspond alors à la distance moyenne entre deux chocs, appelé \imp{libre parcours moyen}. Dans un gaz à température ambiante, il vaut $\ell \simeq \SI{100}{nm}$.


La durée $\tau$ entre deux collisions s’obtient par $\tau \simeq \ell/v$, où $v$ est
la vitesse typique de déplacement des atomes entre deux chocs. Cette vitesse est donnée par analyse dimensionnelle,
\begin{equation}
v \simeq \sqrt{\frac{RT}{M}} \simeq \SI{100}{m.s^{-1}},
\end{equation}
à température ambiante $T \simeq \SI{300}{K}$, pour l’argon, de masse molaire $M=\SI{}{g.mol^{-1}}$.
On obtient pour le coefficient de diffusion,
\begin{equation}
D \simeq \frac{{\ell}^{\,2}}{\tau} \simeq \ell v \simeq \SI{e-5}{m^2.s^{-1}}.
\end{equation}
%où $v^{\star}$ est la vitesse quadratique moyenne des atomes, qui vaut $\sqrt{RT/M}$. Ceci conduit à $\bar{l}\sim \SI{100}{nm}$ et $v^{\star}\sim\SI{100}{m.s^{-1}}$ d'où environ $D\sim\SI{e-5}{m^2.s^{-1}}$,
Ce modèle de marche aléatoire donne un bon ordre de grandeur des coefficients de diffusion des gaz.
% \newpage

\end{document}
